% !PRINTED 451
\weberpage{451}
Nun liegt der Punkt $\pi_0$ auf den drei Hauptaxen der Substitutionen $2^{\mathrm{ten}}$ Grades $\omega, \omega\chi, \omega\chi^2$. Giebt es noch eine weitere Substitution $2^{\mathrm{ten}}$ Grades, durch die $\pi_0$ ungeändert bleibt, so muss $\pi_0$ auf der Hauptaxe dieser Substitution liegen.

Setzen wir aber in den Gleichungen der Hauptaxen [\S.\,118, (7)] $x_1 = x_2 = x_3$, so ist jede dieser drei Gleichungen nur für einen Werth von $r$ befriedigt; also kann $\pi_0$ nicht auf mehr als dreien der Hauptaxen liegen, und $\pi_0$ ist sechszählig.

Daraus der Satz:

7. \gesperrt{Es giebt 42 sechszählige Pole, die alle aus einem von ihnen durch Substitutionen der Gruppe ableitbar sind. Durch jeden dieser Pole gehen drei Hauptaxen der Gruppe.}

Der Pol $\pi_0$ liegt auf keiner der Linien $T_{s,r}$. Dies lässt sich leicht zeigen, wenn man in den Ausdrücken (4) $x_1 = x_2 = x_3$, und dann für $\alpha^2, \beta^2, \gamma^2$ ihre Ausdrücke durch $\varepsilon$ setzt. Man sieht dann leicht (mit Anwendung der Irreducibilität der Kreistheilungsgleichung), dass von diesen Ausdrücken keiner verschwindet.

\S.\,120.

\gesperrt{Die Configuration der Gruppe $G_{168}$.}

Die Gesammtheit der geraden Linien und Punkte, die wir in den vorangegangenen Paragraphen betrachtet haben, wollen wir die \gesperrt{Configuration der Gruppe} $G_{168}$ nennen. Das Wort hat hier dieselbe Bedeutung, in der es in neuerer Zeit in der Geometrie gebraucht wird\footnote{Der Ausdruck ist von \gesperrt{Reye} eingeführt: \glqq{}Geometrie der Lage.\grqq{} Bd. I, 2. Aufl. (1876).}.

Wir beschreiben diese Configuration im Folgenden etwas näher, ohne etwas Neues hinzuzufügen, nur um die Sätze der letzten Paragraphen anschaulicher und übersichtlicher hervortreten zu lassen.

Wir haben in dieser Configuration:
\[
\begin{aligned}
&\text{21 Linien } A \text{ (die Hauptaxen)}, \\
&\text{21 Punkte } a \text{ (die achtzähligen Pole)}, \\
&\text{28 Linien } T, \\
&\text{28 Punkte } t \text{ (die sechszähligen Pole)}.
\end{aligned}
\]
